% !TEX root =./ECE444_Assessment_main.tex
%
% Module 1 assessment content. One question block per learning objective.
% LOs 1.1, 1.2, 1.3, 1.5, 1.6 are stubbed; LO 1.4 is seeded with the
% amplifier / protection-diode mismatch problem.
%
% NOTE on exam-class layout: \fullwidth (a full-width, un-indented banner)
% must NOT appear before the first \question inside a questions list, or the
% class throws "missing \item". So the FIRST LO banner sits just before
% \begin{questions}; every later banner sits between questions, which is legal.

%%%%%%%%%%%%%%%%%%%%%%%%%%%%%%%%%%%%%%%%%
% LO 1.1
%%%%%%%%%%%%%%%%%%%%%%%%%%%%%%%%%%%%%%%%%
\fullwidth{\textbf{Learning Objective 1.1: } I can define what an antenna is, articulate the reciprocity principle, and describe the role of an antenna in a wireless system.}

\begin{questions}

\question Answer each in one or two sentences.
	\begin{parts}
	\part Define an antenna as a transducer: what does it convert, and in which two
	directions does it operate?
		\begin{solution}[1.1in]
		An antenna is a transducer between a \emph{guided} wave (currents/voltages on a
		transmission line) and a \emph{radiated} free-space wave. It operates both ways:
		transmit (guided $\rightarrow$ radiated) and receive (radiated $\rightarrow$ guided).
		\end{solution}
	\part State the reciprocity principle for antennas and give one measurable consequence.
		\begin{solution}[1.3in]
		Reciprocity: an antenna's transmit and receive characteristics are identical ---
		same radiation pattern, gain, and polarization in either mode. Consequence: a
		pattern may be measured in whichever mode is convenient (e.g.\ on receive) and it
		applies equally on transmit.
		\end{solution}
	\part In a wireless link, name the antenna's role at each end and one system quantity
	it sets.
		\begin{solution}[1.3in]
		Transmit end: converts amplifier power into a directed radiated wave --- sets
		$\text{EIRP} = \myp{t} G_t$. Receive end: collects incident power into the receiver
		--- sets the captured power through its effective aperture $A_e$.
		\end{solution}
	\end{parts}

\newpage
%%%%%%%%%%%%%%%%%%%%%%%%%%%%%%%%%%%%%%%%%
% LO 1.2
%%%%%%%%%%%%%%%%%%%%%%%%%%%%%%%%%%%%%%%%%
\fullwidth{\textbf{Learning Objective 1.2: } I can define and calculate fundamental antenna properties: gain, directivity, effective aperture, beamwidth, and sidelobe level.}

\question A pencil-beam antenna has directivity $25\dbi$ and radiation efficiency
$90\%$, operating at $6\ghz$.

\begin{parts}

	\begin{minipage}[t]{0.62\linewidth}
		\part Find the gain in dBi.
		\begin{solution}[1.1in]
		\eq{ G_{\text{dBi}} = 25\db + 10\,\mylog{0.9} = 25 - 0.46 \approx 24.5\dbi }
		\end{solution}
	\end{minipage}%
	\begin{minipage}[t]{0.33\linewidth}
		\raggedright
		\vspace{12pt}
		$G$ = \ansbox[1.5in]{$\approx 24.5\dbi$}
	\end{minipage}

	\begin{minipage}[t]{0.62\linewidth}
		\part Find the effective aperture at 6 GHz.
		\begin{solution}[1.7in]
		\eq{
			\lambda &= \dfrac{c}{f} = \dfrac{3\ee{8}}{6\ee{9}} = 0.05\m \\
			A_e &= \dfrac{G \lambda^{2}}{4\pi} = \dfrac{(284)(0.05)^{2}}{4\pi}
			     \approx 0.057\ \text{m}^{2}
		}
		\end{solution}
	\end{minipage}%
	\begin{minipage}[t]{0.33\linewidth}
		\raggedright
		\vspace{12pt}
		$A_e$ = \ansbox[1.5in]{$\approx 0.057\ \text{m}^{2}$}
	\end{minipage}

	\begin{minipage}[t]{0.62\linewidth}
		\part The measured pattern's first sidelobe is 13 dB below the main-beam peak.
		State the sidelobe level and the aperture illumination it implies.
		\begin{solution}[1.1in]
		$\text{SLL} = -13\db$. A first sidelobe near $-13.3\db$ is the signature of a
		\emph{uniformly} illuminated aperture.
		\end{solution}
	\end{minipage}%
	\begin{minipage}[t]{0.33\linewidth}
		\raggedright
		\vspace{12pt}
		SLL = \ansbox[1.5in]{$-13\db$, uniform}
	\end{minipage}

	\begin{minipage}[t]{0.62\linewidth}
		\part The measured half-power beamwidth is $11\mydeg$ (symmetric beam). Estimate
		the directivity and compare with the spec.
		\begin{solution}[1.5in]
		\eq{ D \approx \dfrac{41253}{\theta_{\text{HP}}^{2}} = \dfrac{41253}{(11)^{2}}
		     \approx 341 \;\longrightarrow\; 10\,\mylog{341} \approx 25.3\dbi }
		Consistent with the $25\dbi$ spec.
		\end{solution}
	\end{minipage}%
	\begin{minipage}[t]{0.33\linewidth}
		\raggedright
		\vspace{12pt}
		$D$ = \ansbox[1.5in]{$\approx 25.3\dbi$}
	\end{minipage}

\end{parts}

\newpage
%%%%%%%%%%%%%%%%%%%%%%%%%%%%%%%%%%%%%%%%%
% LO 1.3
%%%%%%%%%%%%%%%%%%%%%%%%%%%%%%%%%%%%%%%%%
\fullwidth{\textbf{Learning Objective 1.3: } I can determine the polarization of an antenna and describe the bandwidth characteristics of common antenna types.}

\question A wave propagating in $+\zhat$ is
$\vec{E} = \xhat\, 4\cosp{\omega t - kz} + \yhat\, 4\cosp{\omega t - kz - 90\mydeg}$.
The receiving antenna is specified as VSWR $\le 2$ over $2.40$--$2.50\ghz$.

\begin{parts}

	\begin{minipage}[t]{0.62\linewidth}
		\part Identify the polarization of the wave.
		\begin{solution}[1in]
		Equal amplitudes with $\delta = -90\mydeg$ and $+\zhat$ propagation
		$\rightarrow$ \textbf{right-hand circular (RHCP)}.
		\end{solution}
	\end{minipage}%
	\begin{minipage}[t]{0.33\linewidth}
		\raggedright
		\vspace{12pt}
		Polarization = \ansbox[1.5in]{RHCP}
	\end{minipage}

	\begin{minipage}[t]{0.62\linewidth}
		\part A \emph{linearly} polarized transmitter illuminates this antenna. Find the
		polarization loss factor (PLF) in dB.
		\begin{solution}[1in]
		Linear into circular couples half the power:
		$\text{PLF} = 0.5 = -3\db$.
		\end{solution}
	\end{minipage}%
	\begin{minipage}[t]{0.33\linewidth}
		\raggedright
		\vspace{12pt}
		PLF = \ansbox[1.5in]{$-3\db$}
	\end{minipage}

	\begin{minipage}[t]{0.62\linewidth}
		\part Find the fractional bandwidth of the VSWR $\le 2$ band.
		\begin{solution}[1.1in]
		\eq{ \text{FBW} = \dfrac{2.50 - 2.40}{2.45} \approx 0.041 = 4.1\% }
		\end{solution}
	\end{minipage}%
	\begin{minipage}[t]{0.33\linewidth}
		\raggedright
		\vspace{12pt}
		FBW = \ansbox[1.5in]{$\approx 4.1\%$}
	\end{minipage}

	\begin{minipage}[t]{0.62\linewidth}
		\part Convert VSWR $= 2$ to $\left| \Gamma \right|$ and return loss.
		\begin{solution}[1.4in]
		\eq{
			\left| \Gamma \right| &= \dfrac{2 - 1}{2 + 1} = \dfrac{1}{3} \approx 0.333 \\
			\text{RL} &= -20\,\mylog{\tfrac{1}{3}} \approx 9.5\db
		}
		\end{solution}
	\end{minipage}%
	\begin{minipage}[t]{0.33\linewidth}
		\raggedright
		\vspace{12pt}
		$\left| \Gamma \right|$, RL = \ansbox[1.5in]{$\tfrac{1}{3}$, $9.5\db$}
	\end{minipage}

\end{parts}

\newpage
%%%%%%%%%%%%%%%%%%%%%%%%%%%%%%%%%%%%%%%%%
% LO 1.4
%%%%%%%%%%%%%%%%%%%%%%%%%%%%%%%%%%%%%%%%%
\fullwidth{\textbf{Learning Objective 1.4: } I can calculate input impedance, feed considerations, and the role of baluns in an antenna feed system.}

\question A radar transmitter's power amplifier (PA) delivers a forward power of
$\myp{fwd} = +43\dbm$ toward its antenna. A shunt protection diode at the PA
output can safely dissipate at most $2\watts$ of reflected power before
catastrophic failure. On the bench, the antenna measures a VSWR of $2.0:1$.
Model \emph{all} of the reflected power as landing on the diode.

\begin{parts}

	\begin{minipage}[t]{0.62\linewidth}
		\part Convert the forward power to watts.
		\begin{solution}[1.3in]
		\eq{
			\myp{fwd} = 10^{(43-30)/10}\watts = 10^{1.3}\watts \approx 20\watts
		}
		\end{solution}
	\end{minipage}%
	\begin{minipage}[t]{0.33\linewidth}
		\raggedright
		\vspace{12pt}
		$\myp{fwd}$ = \ansbox[1.5in]{$\approx 20\watts$}
	\end{minipage}

	\begin{minipage}[t]{0.62\linewidth}
		\part From the VSWR, find the reflection-coefficient magnitude
		$\left| \Gamma \right|$.
		\begin{solution}[1.3in]
		\eq{
			\left| \Gamma \right| = \dfrac{\text{VSWR} - 1}{\text{VSWR} + 1}
			= \dfrac{2 - 1}{2 + 1} = \dfrac{1}{3} \approx 0.333
		}
		\end{solution}
	\end{minipage}%
	\begin{minipage}[t]{0.33\linewidth}
		\raggedright
		\vspace{12pt}
		$\left| \Gamma \right|$ = \ansbox[1.5in]{$\tfrac{1}{3} \approx 0.333$}
	\end{minipage}

	\begin{minipage}[t]{0.62\linewidth}
		\part Compute the reflected power reaching the diode.
		\begin{solution}[1.3in]
		\eq{
			\myp{ref} = \left| \Gamma \right|^{2} \myp{fwd}
			= \lp \dfrac{1}{3} \rp^{2} (20\watts)
			= 0.111 (20\watts) \approx 2.22\watts
		}
		\end{solution}
	\end{minipage}%
	\begin{minipage}[t]{0.33\linewidth}
		\raggedright
		\vspace{12pt}
		$\myp{ref}$ = \ansbox[1.5in]{$\approx 2.22\watts$}
	\end{minipage}

	\begin{minipage}[t]{0.62\linewidth}
		\part Is the diode safe at this VSWR? Justify in one line.
		\begin{solution}[1.1in]
		$\myp{ref} \approx 2.22\watts > 2\watts$ rating
		$\longrightarrow$ \textbf{not safe}; the diode is over-driven and fails.
		\end{solution}
	\end{minipage}%
	\begin{minipage}[t]{0.33\linewidth}
		\raggedright
		\vspace{12pt}
		Safe? \ansbox[1.5in]{No, $2.22 > 2\watts$}
	\end{minipage}

	\begin{minipage}[t]{0.62\linewidth}
		\part What is the maximum VSWR this system can tolerate before the diode
		is destroyed?
		\begin{solution}[2in]
		\eq{
			\left| \Gamma \right|_{\max}
				&= \sqrt{\dfrac{\myp{max}}{\myp{fwd}}}
				= \sqrt{\dfrac{2\watts}{20\watts}} = \sqrt{0.1} \approx 0.316 \\
			\text{VSWR}_{\max}
				&= \dfrac{1 + \left| \Gamma \right|_{\max}}
				         {1 - \left| \Gamma \right|_{\max}}
				= \dfrac{1.316}{0.684} \approx 1.93
		}
		\end{solution}
	\end{minipage}%
	\begin{minipage}[t]{0.33\linewidth}
		\raggedright
		\vspace{12pt}
		$\text{VSWR}_{\max}$ = \ansbox[1.5in]{$\approx 1.93$}
	\end{minipage}

\end{parts}

\newpage
%%%%%%%%%%%%%%%%%%%%%%%%%%%%%%%%%%%%%%%%%
% LO 1.5
%%%%%%%%%%%%%%%%%%%%%%%%%%%%%%%%%%%%%%%%%
\fullwidth{\textbf{Learning Objective 1.5: } I can identify and distinguish the reactive near-field, radiating near-field, and far-field regions and calculate the boundaries for a given antenna.}

\question An aperture antenna has largest dimension $D = 0.5\m$ and operates at $3\ghz$.

\begin{parts}

	\begin{minipage}[t]{0.62\linewidth}
		\part Find the wavelength.
		\begin{solution}[0.9in]
		\eq{ \lambda = \dfrac{c}{f} = \dfrac{3\ee{8}}{3\ee{9}} = 0.10\m }
		\end{solution}
	\end{minipage}%
	\begin{minipage}[t]{0.33\linewidth}
		\raggedright
		\vspace{12pt}
		$\lambda$ = \ansbox[1.5in]{$0.10\m$}
	\end{minipage}

	\begin{minipage}[t]{0.62\linewidth}
		\part Find the far-field (Fraunhofer) boundary, $2D^{2}/\lambda$.
		\begin{solution}[1in]
		\eq{ r_{\text{ff}} = \dfrac{2D^{2}}{\lambda} = \dfrac{2(0.5)^{2}}{0.10} = 5\m }
		\end{solution}
	\end{minipage}%
	\begin{minipage}[t]{0.33\linewidth}
		\raggedright
		\vspace{12pt}
		$r_{\text{ff}}$ = \ansbox[1.5in]{$5\m$}
	\end{minipage}

	\begin{minipage}[t]{0.62\linewidth}
		\part Find the outer limit of the reactive near-field, $0.62\sqrt{D^{3}/\lambda}$.
		\begin{solution}[1.2in]
		\eq{ r_{\text{rnf}} = 0.62\sqrt{\dfrac{D^{3}}{\lambda}}
		     = 0.62\sqrt{\dfrac{(0.5)^{3}}{0.10}} = 0.62\sqrt{1.25} \approx 0.69\m }
		\end{solution}
	\end{minipage}%
	\begin{minipage}[t]{0.33\linewidth}
		\raggedright
		\vspace{12pt}
		$r_{\text{rnf}}$ = \ansbox[1.5in]{$\approx 0.69\m$}
	\end{minipage}

	\begin{minipage}[t]{0.62\linewidth}
		\part How far must an antenna range place the AUT for a valid far-field pattern
		measurement, and why?
		\begin{solution}[1in]
		At least the far-field distance, $r \ge 2D^{2}/\lambda = 5\m$, so the phase across
		the aperture is within the $22.5\mydeg$ criterion and the wavefront is effectively
		planar.
		\end{solution}
	\end{minipage}%
	\begin{minipage}[t]{0.33\linewidth}
		\raggedright
		\vspace{12pt}
		$r_{\min}$ = \ansbox[1.5in]{$\ge 5\m$}
	\end{minipage}

\end{parts}

\newpage
%%%%%%%%%%%%%%%%%%%%%%%%%%%%%%%%%%%%%%%%%
% LO 1.6
%%%%%%%%%%%%%%%%%%%%%%%%%%%%%%%%%%%%%%%%%
\fullwidth{\textbf{Learning Objective 1.6: } I can set up and interpret the radiation integrals to derive the far-field pattern of a current distribution.}

\question A thin, $z$-directed current filament carries a \emph{uniform} current $I_0$
over $-L/2 \le z^{\prime} \le L/2$. Its far-field pattern factor is the radiation integral
$N(\theta) = \displaystyle\int I_0\, e^{+j k z^{\prime} \cosp{\theta}}\, dz^{\prime}$.

\begin{parts}

	\part Evaluate the integral $N(\theta)$ in closed form.
		\begin{solution}[1.8in]
		\eq{
			N(\theta) &= I_0 \int_{-L/2}^{L/2} e^{+j k z^{\prime} \cosp{\theta}}\, dz^{\prime}
			         = I_0\, \dfrac{2\sinp{\tfrac{kL}{2}\cosp{\theta}}}{k\cosp{\theta}} \\
			         &= I_0 L \; \dfrac{\sinp{\tfrac{kL}{2}\cosp{\theta}}}
			                          {\tfrac{kL}{2}\cosp{\theta}}
		}
		\end{solution}

	\part Write the normalized far-field pattern $\left| F(\theta) \right|$.
		\begin{solution}[1.1in]
		\eq{ \left| F(\theta) \right|
		     = \left| \dfrac{\sinp{\tfrac{kL}{2}\cosp{\theta}}}
		                    {\tfrac{kL}{2}\cosp{\theta}} \right| }
		\end{solution}

	\part Evaluate $\left| F(\theta) \right|$ at broadside, $\theta = 90\mydeg$, and
	explain the result physically.
		\begin{solution}[1.3in]
		At $\theta = 90\mydeg$, $\cosp{\theta} = 0$, so the argument $\to 0$ and
		$\tfrac{\sin x}{x} \to 1$: thus $\left| F(90\mydeg) \right| = 1$. The uniform line
		source radiates maximally \emph{broadside} (perpendicular) to its axis.
		\end{solution}

\end{parts}

\end{questions}
